\section{Virtualização de Temporizador}

Por definição da arquitetura RISC-V, na ausência de extensões específicas de hardware, os recursos de temporizador — especificamente os registradores \textit{mtime} e \textit{mtimecmp} —
são exclusivos e gerenciados pelo modo \textit{Machine}.

Para contornar essa restrição em ambientes virtualizados, os \glspl{so} convidados --- executados em modo S ou VS --- utilizam a \gls{sbi} para gerenciar eventos temporais.
Dessa forma, não é necessária nenhuma adaptação estrutural no \gls{so} convidado, uma vez que a própria especificação padrão da arquitetura prevê o uso da \gls{sbi} para essa finalidade.

\subsection{Leitura do Temporizador}

A leitura do tempo atual pelo convidado é feita por meio do \gls{csr} \textit{time}, acessado comumente através da pseudo-instrução \textit{rdtime}.

Quando um modo com menor privilégio tenta realizar essa leitura sem que o acesso direto tenha sido liberado, o processador dispara uma \textit{Illegal Instruction}. O hipervisor intercepta essa exceção
calcula o tempo virtualizado correspondente e devolve o valor correto ao registrador de destino do convidado.

\subsection{Programação de Temporizador}

Como o convidado não possui acesso direto de escrita ao registrador \textit{mtimecmp}, o agendamento de interrupções é realizado por chamadas \gls{sbi} --- extensão \textit{TIME}.

Quando o \gls{so} convidado precisa agendar um evento futuro, ele realiza uma chamada de sistema solicitando o registro do próximo limite temporal. O firmware
recebe a requisição e ao atingir o tempo programado, injeta a interrupção sinalizando o bit \gls{stip} do registrador \gls{mip} da CPU virtual correspondente.
